\documentclass[aspectratio=169,11pt]{beamer}

\usetheme{Madrid}
\usecolortheme{default}
\usefonttheme{professionalfonts}
\setbeamertemplate{navigation symbols}{}
\setbeamertemplate{footline}[frame number]

\usepackage{amsmath, amssymb, booktabs}

\newcommand{\E}{\mathbb{E}}
\newcommand{\1}{\mathbf{1}}
\newcommand{\plim}{\operatorname{plim}}

\title{LLM Workflows For Applied Research}
\subtitle{Empirical Methods --- Week 14}
\author{Teaching deck draft}
\date{}

\begin{document}

\begin{frame}
  \titlepage
\end{frame}

\begin{frame}{Week 14 question}
\begin{itemize}
    \item How can language models support applied research without weakening reproducibility?
    \item Lecture 14 is about LLMs as research infrastructure, not another text classifier.
    \item The central object is the workflow: prompt, model, retrieval, schema, benchmark, review, and audit trail.
    \item The standard remains applied economics: define the object, validate the measure, preserve the replication path.
\end{itemize}
\end{frame}

\begin{frame}{Five roles for LLMs}
\small
\begin{tabular}{p{0.24\linewidth} p{0.31\linewidth} p{0.34\linewidth}}
\toprule
Role & Output & Main evidentiary question \\
\midrule
Research assistant & code, summaries, templates & Did the researcher verify the final artifact? \\
Coder / extractor & labels, fields, scores & Does it measure the intended economic object? \\
Synthetic annotator & scalable labels & Does it agree with a benchmark and where does it fail? \\
Synthetic subject & choices, survey answers & Does it reproduce human patterns on benchmark tasks? \\
Simulation engine & actions, transitions, narratives & Does the environment generate disciplined comparative statics? \\
\bottomrule
\end{tabular}

\medskip
The role determines the validation standard.
\end{frame}

\begin{frame}{LLMs as research assistants}
\begin{itemize}
    \item High-value uses: code drafting, debugging, cleaning plans, document triage, extraction templates, literature maps.
    \item Strongest when the final published object is still researcher-controlled.
    \item Agentic workflows add tool calls, file edits, and intermediate artifacts.
    \item Every multi-step assistant workflow needs checkpoints and logs.
\end{itemize}

\medskip
Assistant use is not automatically a measurement problem, but it is still a reproducibility problem.
\end{frame}

\begin{frame}{LLMs as coders and extractors}
\small
\[
\widehat z_i = g_{\theta}(x_i,p_i,r_i)
\]

\[
Y_i=\alpha+\beta \widehat z_i+u_i
\]

\begin{itemize}
    \item \(x_i\): raw document, posting, contract, response, or case note.
    \item \(p_i\): prompt and schema.
    \item \(r_i\): retrieval context, examples, rubrics, or taxonomies.
    \item \(\theta\): model version, decoding, tools, and post-processing.
\end{itemize}

\medskip
If \(\widehat z_i\) enters the economics, the workflow is a measurement instrument.
\end{frame}

\begin{frame}{Synthetic coders and annotators}
\begin{itemize}
    \item Useful when human coding is costly and the construct is well specified.
    \item Not enough to show that labels look plausible.
    \item Need human benchmark labels, disagreement audits, and subgroup diagnostics.
    \item Route ambiguous, high-stakes, and low-confidence cases to review.
\end{itemize}

\medskip
A good synthetic-label pipeline is usually human-in-the-loop, not human-free.
\end{frame}

\begin{frame}{LLMs as experimental subjects}
\small
\[
a_i=s_{\theta}(e_i,p_i,c_i)
\]

\begin{itemize}
    \item \(e_i\): vignette, game, policy scenario, or experimental environment.
    \item \(p_i\): persona, instructions, ordering, and prompt frame.
    \item \(c_i\): supplied context or memory.
    \item \(a_i\): simulated choice, response, message, or action.
\end{itemize}

\medskip
Synthetic subjects can pilot designs and explore mechanisms. They do not estimate human behavior unless benchmarked against human behavior.
\end{frame}

\begin{frame}{LLMs as simulation engines}
\small
\[
x_{t+1}=f_{\theta}(x_t,a_t,\varepsilon_t)
\]

\begin{tabular}{p{0.28\linewidth} p{0.62\linewidth}}
\toprule
Design choice & What to make explicit \\
\midrule
State & variables, memory, constraints, histories \\
Action & feasible choices and timing \\
Transition & rule-based versus model-generated updates \\
Shock process & exogenous variation and random seeds \\
Benchmark & empirical regularities the simulation must reproduce \\
\bottomrule
\end{tabular}

\medskip
Fluent narratives are not validation.
\end{frame}

\begin{frame}{Workflow risk decomposition}
\small
\[
R_{\text{workflow}}
=R_{\text{prompt}}
+R_{\text{model}}
+R_{\text{retrieval}}
+R_{\text{benchmark}}
+R_{\text{transport}}
+R_{\text{privacy}}
\]

\begin{itemize}
    \item Prompt risk: wording changes labels or fields.
    \item Model risk: versions, providers, or decoding alter outputs.
    \item Retrieval risk: wrong chunks, stale corpus, missing context.
    \item Benchmark risk: small, easy, biased, or nonportable gold labels.
    \item Transport risk: source prompt does not travel to a new domain.
    \item Privacy risk: raw inputs cannot safely leave the data environment.
\end{itemize}
\end{frame}

\begin{frame}{Replication, audit trails, and privacy}
\small
\begin{tabular}{p{0.28\linewidth} p{0.62\linewidth}}
\toprule
Artifact & Minimum record \\
\midrule
Prompt & system prompt, user template, examples, schema \\
Model & provider, model, version/date, decoding settings \\
Retrieval & corpus version, chunking, retrieved passages, settings \\
Output & raw completion, parsed fields, validation errors \\
Review & reviewer, adjudication, reason code, final label \\
Evaluation & benchmark set, metrics, sensitivity scripts \\
Privacy & release limits, redaction, secure review protocol \\
\bottomrule
\end{tabular}

\medskip
Each published value should be traceable from raw input to final dataset.
\end{frame}

\begin{frame}{Publication-ready LLM measurement}
\begin{itemize}
    \item Define the economic object before choosing the model.
    \item Freeze a schema, prompt, retrieval corpus, and post-processing rule.
    \item Build or preserve an independent benchmark set.
    \item Report accuracy, calibration, subgroup performance, and prompt sensitivity.
    \item Audit false positives, false negatives, unsupported completions, and influential cases.
    \item Explain how remaining errors change descriptive, causal, structural, or policy interpretation.
\end{itemize}
\end{frame}

\begin{frame}{Research Lab design}
\begin{columns}[T,onlytextwidth]
\begin{column}{0.32\linewidth}
\textbf{Reproduce}
\begin{itemize}
    \item synthetic labor documents
    \item prompt template
    \item structured schema
    \item prompt log
    \item benchmark metrics
\end{itemize}
\end{column}
\begin{column}{0.32\linewidth}
\textbf{Diagnose}
\begin{itemize}
    \item prompt sensitivity
    \item subgroup stability
    \item unsupported outputs
    \item human-review queue
    \item downstream slope
\end{itemize}
\end{column}
\begin{column}{0.32\linewidth}
\textbf{Transfer}
\begin{itemize}
    \item workforce case notes
    \item source prompt
    \item adapted prompt
    \item transport metrics
    \item redesign memo
\end{itemize}
\end{column}
\end{columns}
\end{frame}

\begin{frame}{Takeaways}
\begin{itemize}
    \item LLMs are useful research infrastructure when their role is explicit.
    \item A prompt/model workflow can be a measurement instrument.
    \item Synthetic subjects are best treated as pilots or benchmarked objects of study.
    \item Simulation engines need transparent state, action, transition, and validation rules.
    \item Replication requires logs, outputs, schemas, benchmarks, review records, and privacy documentation.
    \item LLMs lower some costs; they do not lower the standard of evidence.
\end{itemize}
\end{frame}

\end{document}
